% Section 13.4: General Infrared Divergences.
% Source: physical PDF pp. 571--576 (printed pp. 548--553).
% Starts with the section heading below Eq. (13.3.12) on p. 571 and ends with
% Eq. (13.4.13), its following sentence, and the attached note on p. 576.
% Equations: (13.4.1)--(13.4.13). Figures: none. Footnotes: three.
% Status: source-reviewed against rendered pages; no unresolved uncertainties.

\subsection{General Infrared Divergences}

The infrared divergence due to soft photons that we have been considering up
to now in this chapter is just one example of a variety of infrared divergences
that are encountered in various physical theories. Another example is
provided by quantum electrodynamics with massless charged particles. Here even
after the cancellation of infrared divergences due to soft photons, we find a
logarithmic divergence in the exponent $A$ in \eqref{eq:13.3.11}.
According to Eqs. \eqref{eq:13.2.11} and
\eqref{eq:13.2.7}, for a process in which all charged particles are
electrons, in the limit $m_e\to0$ the exponent goes as
\[
A\to-\frac{1}{4\pi^2}\sum_n e_n^2
-\frac{1}{4\pi^2}\sum_{n\ne m}e_ne_m\eta_n\eta_m
\ln\left(\frac{2|p_n\cdot p_m|}{m_e^2}\right)
\to-\frac{\ln m_e}{2\pi^2}\sum_n e_n^2.
\]
(In the last step we have used the charge conservation condition
$\sum_ne_n\eta_n=0$.) The infrared divergence in this formula arises from
soft photons that are emitted in a direction \emph{parallel} to the momentum
of one of the `hard' electrons in the initial or final state, but it occurs
also even if the photon like the electron is not soft, because the propagator
denominator $(p_n\pm q)^2$ vanishes for $p_n^2=q^2=0$ if $\mathbf p_n$ is
parallel to $\mathbf q$. To be a little more specific, for $p_n^2=q^2=0$ the
integral of this factor\footnote{This factor is not squared, because the
divergence occurs only in the interference between this term in the $S$-matrix
element and terms in which the photon is emitted from some other charged
particle line $m\ne n$. For $m=n$ the integral \eqref{eq:13.2.8} is
proportional to $m_n^2$.} over photon directions takes the form
\[
\int d^2\hat{\mathbf q}\,(p\pm q)^{-2}
=\frac{2\pi}{\mathbf p^2\mathbf q^2}\int_0^\pi
\frac{\sin\theta\,d\theta}{1-\cos\theta},
\]
where $\theta$ is the angle between the momenta of the photon and the charged
particle. This integral diverges logarithmically at $\theta=0$.

Of course, in the real world there are no massless electrically charged
particles, but in reactions in which the typical value $E^2$ of the scalar
products $|p_n\cdot p_m|$ is much larger than $m_e^2$, it is of interest to
identify the places where large $\ln(m_e/E)$ factors appear. The dominant
radiative correction in this case is often given by the term
$-\ln(m_e/E)\sum_n e_n^2/2\pi^2$ in $A$. More importantly, in quantum
chromodynamics there are massless particles, the gluons, that carry a
conserved quantum number known as color that is analogous to electric charge,
so that infrared divergences arise from the emission of parallel hard gluons
from hard gluons or other hard colored particles in the initial or final
states.

These infrared divergences are not, in general, eliminated by summing over
suitable sets of final states. However, Lee and
Nauenberg~\hyperref[ref:13.6]{[6]} have pointed out that the infrared
divergences can be made to cancel if we not only sum over suitable final
states, but also assume a certain probabilistic distribution of \emph{initial}
states. What follows is a modified version of their argument, which will
immediately make clear why in the case of electrodynamics with massive
electrons it was sufficient to sum over final states.

For these purposes it is convenient to return to `old-fashioned' perturbation
theory, in which the $S$-matrix is given by Eq. \eqref{eq:3.2.7} and
Eq. \eqref{eq:3.5.3} as
\begin{equation}
S_{ba}=\delta(b-a)-2i\pi\delta(E_a-E_b)T_{ba},
\tag{13.4.1}\label{eq:13.4.1}
\end{equation}
where
\begin{equation}
T_{ba}=(H_{\mathrm{int}})_{ba}
+\sum_{\nu=1}^{\infty}\int dc_1\cdots dc_\nu\,
\frac{(H_{\mathrm{int}})_{bc_\nu}
(H_{\mathrm{int}})_{c_\nu c_{\nu-1}}\cdots
(H_{\mathrm{int}})_{c_1a}}
{(E_a-E_{c_1}+i\epsilon)\cdots(E_a-E_{c_\nu}+i\epsilon)}.
\tag{13.4.2}\label{eq:13.4.2}
\end{equation}
(The integrals over $c_1\cdots c_\nu$ should be understood to include sums
over the spins and types of particles in these states as well as integrals
over the three-momenta of these particles.) Infrared divergences arise from
(and only from) the vanishing of one or more of the energy denominators in
this expression.

However, not all vanishing energy denominators give rise to infrared
divergences. A general intermediate state $c$ may have $E_c=E_a$, but usually
this is just one point in the interior of the range of integration, and the
integral over this range is rendered convergent by the prescription implied by
the $i\epsilon$ in the denominator. In order for an intermediate state $c$ to
produce an infrared divergence, it is necessary that the energy $E_c=E_a$ be
reached at the \emph{endpoint} of the range of integration. This happens for
instance if the first intermediate state $c_1$ in \eqref{eq:13.4.2}
consists of the particles in the initial state $a$, with any of the massless
particles in this state replaced with \emph{jets}, consisting of any number of
nearly parallel massless particles with a total momentum equal to that of the
particle the jet replaces. In this case, the endpoint at which $E_{c_1}=E_a$
is the point in momentum space at which all of the massless particles in each
jet are parallel. More generally, we can have any number of the massless
particles in $a$ replaced with jets of nearly parallel massless particles,
plus any number of additional soft massless particles. The set of all such
states will be called $D(a)$. (To be precise, we need to introduce a small
angle $\Theta$ as well as a small energy $\Lambda$ to define what we mean by
`nearly parallel' and `soft'. We will not bother to show the dependence of the
set $D(a)$ on $\Theta$ and $\Lambda$.) The states in $D(a)$ are `dangerous', in
the sense that the vanishing of the energy denominator $E_a-E_{c_1}$ at the
endpoint can introduce an infrared divergence; the endpoint at which
$E_{c_1}=E_a$ is the point at which all massless particles in each jet are
parallel, and all soft massless particles have zero energy.

Furthermore, if $c_1,\ldots,c_n$ are each in the set $D(a)$, then an
intermediate state $c_{n+1}$ in $D(a)$ is also dangerous in the same sense. On
the other hand, if some intermediate state $c_m$ is \emph{not} in $D(a)$, then
a later state $c_k$ with $k>m$ would not be dangerous even if it belonged to
the set $D(a)$, because the configuration of hard particles or jets with
three-momenta equal to those of particles in the state $a$ would be just an
ordinary point inside the range of integration. In exactly the same way, we
may define a set of states $D(b)$ in which one or more of the massless
particles in the state $b$ are replaced with jets of nearly parallel massless
particles, each having the same total three-momentum as the particle it
replaces, and we add any number of soft massless particles. An intermediate
state $c_m$ is dangerous if it belongs to the set $D(b)$ and if the later
states $c_k$ with $k>m$ all belong to $D(b)$.

To isolate the effects of these dangerous states we rewrite
\eqref{eq:13.4.2} in the form
\begin{equation}
T_{ba}=(H_{\mathrm{int}})_{ba}
+\sum_{\nu=1}^{\infty}
\left(H_{\mathrm{int}}
\left[\frac{\mathcal P_a+\mathcal P_b+\mathcal P_{\notin a,b}}
{E_a-H_0+i\epsilon}H_{\mathrm{int}}\right]^\nu\right)_{ba},
\tag{13.4.3}\label{eq:13.4.3}
\end{equation}
where $\mathcal P_a$, $\mathcal P_b$, and $\mathcal P_{\notin a,b}$ are the
projection operators respectively on $D(a)$, $D(b)$, and on all other states.
(It is assumed here that none of the charged particles in $b$ have momenta
close to that of some charged particle in $a$, so that the sets $D(a)$ and
$D(b)$ do not overlap.) Now, for $\Lambda\to0$ and $\Theta\to0$, the dangerous
intermediate states occupy so little phase space that they may be neglected
wherever they do not lead to infrared divergences. The power series
\eqref{eq:13.4.3} therefore becomes
\begin{align}
T_{ba}={}&\sum_{r=0}^{\infty}\sum_{s=0}^{\infty}
\sum_{\nu=0}^{\infty}
\Biggl(
\left[H_{\mathrm{int}}\frac{\mathcal P_b}{E_a-H_0+i\epsilon}\right]^r
H_{\mathrm{int}}
\left[\frac{\mathcal P_{\notin a,b}}{E_a-H_0+i\epsilon}
H_{\mathrm{int}}\right]^\nu\notag\\
&\hspace{33mm}\cdot
\left[\frac{\mathcal P_a}{E_a-H_0+i\epsilon}
H_{\mathrm{int}}\right]^s
\Biggr)_{ba}.
\tag{13.4.4}\label{eq:13.4.4}
\end{align}
This would be exact if all of the projection operators
$\mathcal P_{\notin a,b}$ between the leftmost and the rightmost were replaced
with $\mathcal P_a+\mathcal P_b+\mathcal P_{\notin a,b}$, and if
$\mathcal P_b$ and $\mathcal P_a$ on the left and right were replaced with
$\mathcal P_b+\mathcal P_a$, but as remarked above this would have a negligible
effect on the final result when $\Lambda$ and $\Theta$ are sufficiently small.

Eq. \eqref{eq:13.4.4} may be written in a more compact form:
\begin{equation}
T_{ba}=\left\{(\Omega_b^-)^\dagger T_S\Omega_a^+\right\}_{ba},
\tag{13.4.5}\label{eq:13.4.5}
\end{equation}
where, for future use, we define $\Omega_\alpha^+$ and $\Omega_\beta^-$ for
general states $\alpha$ and $\beta$ as:
\begin{equation}
(\Omega_\alpha^+)_{ca}\equiv
\sum_{r=0}^{\infty}
\left(\left[\frac{\mathcal P_\alpha}{E_\alpha-H_0+i\epsilon}
H_{\mathrm{int}}\right]^r\right)_{ca},
\tag{13.4.6}\label{eq:13.4.6}
\end{equation}
\begin{equation}
(\Omega_\beta^-)_{db}\equiv
\sum_{r=0}^{\infty}
\left(\left[\frac{\mathcal P_\beta}{E_\beta-H_0-i\epsilon}
H_{\mathrm{int}}\right]^r\right)_{db},
\tag{13.4.7}\label{eq:13.4.7}
\end{equation}
and $T_S$ is the `safe' operator\footnote{In $(\Omega_b^-)_{db}$ we are using
the fact that $T_{ba}$ is calculated with $E_b=E_a$, and in $(T_S)_{dc}$ we
are using the fact that the projection operators $\mathcal P_a$ make
$(\Omega_a^+)_{ca}$ vanish unless $E_c$ is very close to $E_a$. Also, the
factors $(\Omega_b^-)^\dagger$ and $\Omega_a^+$ in
\eqref{eq:13.4.5} make
$\mathcal P_{\notin c,d}=\mathcal P_{\notin a,b}$.}
\begin{equation}
(T_S)_{dc}\equiv\sum_{\nu=0}^{\infty}
\left(H_{\mathrm{int}}
\left[\frac{\mathcal P_{\notin c,d}}{E_c-H_0+i\epsilon}
H_{\mathrm{int}}\right]^\nu\right)_{dc}.
\tag{13.4.8}\label{eq:13.4.8}
\end{equation}
All of the infrared divergences have now been isolated in the two operator
factors $\Omega_\beta^-$ and $\Omega_\alpha^+$.

To eliminate these infrared divergences it is now only necessary to note that
if it were \emph{not} for the projection operators on the dangerous states,
the operators $\Omega_\beta^-$ and $\Omega_\alpha^+$ would be just the unitary
operators that according to Eq. \eqref{eq:3.1.16} convert free-particle states
into `out' or `in' states, respectively. These operators are therefore
\emph{unitary if confined to the subspaces $D(\beta)$ and $D(\alpha)$ of
states that would be dangerous for some given final state $\beta$ and some
given initial state $\alpha$}. That is, for general $\alpha$ and $\beta$
\begin{equation}
\Omega_\beta^-\mathcal P_\beta(\Omega_\beta^-)^\dagger
=\mathcal P_\beta,
\tag{13.4.9}\label{eq:13.4.9}
\end{equation}
\begin{equation}
\Omega_\alpha^+\mathcal P_\alpha(\Omega_\alpha^+)^\dagger
=\mathcal P_\alpha.
\tag{13.4.10}\label{eq:13.4.10}
\end{equation}

\emph{The transition rate is therefore free of infrared divergences if
summed over the subspaces of states that would be dangerous for any given
final and initial states $\beta$ and $\alpha$:}
\begin{align}
\sum_{a\in D(\alpha)}\sum_{b\in D(\beta)}|T_{ba}|^2
={}&\operatorname{Tr}\left\{
\Omega_\beta^-\mathcal P_\beta(\Omega_\beta^-)^\dagger
T_S\Omega_\alpha^+\mathcal P_\alpha(\Omega_\alpha^+)^\dagger T_S^\dagger
\right\}\notag\\
={}&\operatorname{Tr}\left\{\mathcal P_\beta T_S
\mathcal P_\alpha T_S^\dagger\right\}
=\sum_{a\in D(\alpha)}\sum_{b\in D(\beta)}|(T_S)_{ba}|^2.
\tag{13.4.11}\label{eq:13.4.11}
\end{align}

In order to be satisfied that this really does solve the general problem of
infrared divergences, it is necessary to argue that it is only sums like that
in \eqref{eq:13.4.11} that are experimentally measurable. It is
plausible that we should have to sum over dangerous final states in order to
have a measurable transition rate, since it is not possible experimentally to
distinguish an outgoing charged (or colored) massless particle from a jet of
massless particles with nearly parallel momenta and the same total
energy~\hyperref[ref:13.7]{[7]}, together with an arbitrary number of very
soft quanta, all with the same total charge (or color). The sum over initial
states is more problematic. Presumably one may argue that truly massless
particles are always produced as jets accompanied by an ensemble of soft
quanta that is uniform within some volume of momentum space. However, to the
best of my knowledge no one has given a complete demonstration that the sums
of transition rates that are free of infrared divergences are the only ones
that are experimentally measurable.

This problem does not arise in quantum electrodynamics (with massive charged
particles), where as we have seen it is only necessary to sum over final states
in order to eliminate infrared divergences. The reason for this difference can
be traced to the fact that in electrodynamics the states $a,b,c,\ldots$ are
direct products of states (labelled with Greek letters) with fixed numbers of
charged particles and hard photons, times states containing only soft photons
having energy less than some small quantity $\Lambda$. Then for a reaction in
which some set of soft photons $f$ is produced in a reaction $\alpha\to\beta$
among charged particles and hard photons, \eqref{eq:13.4.5}
simplifies to
\begin{equation}
T_{\beta f,\alpha}
=\langle f|(\Omega^-(\beta))^\dagger\Omega^+(\alpha)|\Omega\rangle
(T_S)_{\beta\alpha},
\tag{13.4.12}\label{eq:13.4.12}
\end{equation}
where $|\Omega\rangle$ denotes the soft photon vacuum, and $\Omega^\pm$ are
calculated as before, but in the reduced Hilbert space consisting only of soft
photons, and with the interactions of these photons taken as the interaction
Hamiltonian with all charged particles in the fixed states indicated by the
arguments $\beta$ or $\alpha$. Just as before, these operators are unitary in
the `dangerous' Hilbert space $\mathcal D$ of soft photons,
so\footnote{The reason that we are now not encountering any factor
$(E/\Lambda)^A$ like that in \eqref{eq:13.3.11} is that we are here
identifying the maximum energy $E$ of the real soft photon states over which we
are summing with the maximum energy $\Lambda$ of the `dangerous' soft photon
states over which we sum in calculating $\Omega^\pm$.}
\begin{align}
\sum_{f\in\mathcal D}|T_{\beta f,\alpha}|^2
={}&|(T_S)_{\beta\alpha}|^2
\langle\Omega|(\Omega^+(\alpha))^\dagger\Omega^-(\beta)
(\Omega^-(\beta))^\dagger\Omega^+(\alpha)|\Omega\rangle\notag\\
={}&|(T_S)_{\beta\alpha}|^2
\langle\Omega|(\Omega^+(\alpha))^\dagger
\Omega^+(\alpha)|\Omega\rangle
=|(T_S)_{\beta\alpha}|^2
\tag{13.4.13}\label{eq:13.4.13}
\end{align}
without having to sum over initial states.
